\section{Method}
\label{sec:method}

Fig.~\ref{fig:pipeline} shows the assembled system. The safety decision is made
outside the frozen policy and enforced upon it, and the labor is divided so that
every neural component answers a local, checkable question while a fixed
symbolic layer makes every decision that matters. We describe the four layers in
the order in which the runtime executes them.

\subsection{Perception}
\label{sec:percep}

\textbf{What is observed.} At \textit{Tier-Percep} the system reads only the
instruction, camera observations, proprioception, and the symbolic list of
object names present in the scene. It never reads
\texttt{obs[\{name\}\_pos]}. Three camera streams are available, and they play
three distinct roles rather than being pooled into one estimate. The
\texttt{agentview} camera is the trusted primary and is the only view used by
default. The overhead \texttt{birdview} camera is a secondary that may refine
the primary but is never permitted to override it. The wrist-mounted
\texttt{robot0\_eye\_in\_hand} camera is consumed by the policy itself as one of
its two image inputs, and the safety layer draws on it only for re-detecting
moving hazards, where at close range the wrist sees a human hand large and
unoccluded in exactly the configurations in which the agentview detector is
weakest.

\textbf{Object inventory.} The candidate set $\mathcal{X}$ is the symbolic name
list, which supplies identity strings but no geometry. Each name is converted
into an open-vocabulary text query by stripping instance suffixes and
underscores, and is passed to GroundingDINO~\cite{gdino} at a box threshold of
0.30. Localization is therefore open-vocabulary in the sense that matters here,
namely that no fixed detector vocabulary, no benchmark-specific class list, and
no per-object model are involved. GroundingDINO runs in a persistent subprocess,
so the detector environment stays out of the evaluation process, and it is
queried once per object per camera at episode start, which makes its cost
negligible relative to policy inference.

\textbf{Depth back-projection.} A detector box is a rectangle, and a rectangle
drawn on a table contains a substantial fraction of table and background pixels,
whose depths would bias a naive centroid. The estimator therefore anchors on the
median depth inside the box, retains the depth-continuous connected component
around that anchor, and back-projects the component's centroid pixel at its
median depth, $p_{\mathrm{cam}}=z\,K^{-1}[u,v,1]^{\top}$ followed by the camera
extrinsic transform, using intrinsics and extrinsics read from the live
simulator. Candidate boxes are tried in descending detector confidence and the
first whose back-projection lands inside the workspace is accepted, which
removes the dominant open-vocabulary failure mode, namely confident detections
on walls and background structure. Half-extents are obtained from the same
masked point cloud as the per-axis fifth to ninety-fifth percentile spread,
clipped to $[0.02,0.35]$\,m. Visible-surface bias underestimates the occluded
axis, and that residual is absorbed by the enforcement pad rather than
corrected.

\textbf{Multi-camera fusion.} Views are fused by mean only when they agree.
Estimates from the primary and secondary cameras are averaged when their maximum
pairwise separation is within $\tau=0.06$\,m, and on disagreement the primary
view wins alone. This asymmetry is not a stylistic preference. Blind mean fusion
imports secondary-view failures wholesale, and in one measured arm a single
mislocalized birdview detection collapsed two obstacle-avoidance task cells to
near-zero success while leaving the intended target cell essentially unchanged.
Under the consistency gate a secondary view can only refine an estimate, never
replace it. Median localization error is 0.067\,m on SafeLIBERO and
0.036--0.038\,m on LIBERO-Safety.

\textbf{Failure behavior.} An object whose estimate fails is simply absent from
the geometric layer, so it can neither be elected as a hazard nor grant a
sanctioned corridor. Failures are logged per episode rather than silently
imputed. Dynamic hazards are handled by re-perception rather than by prediction:
on the human-hazard suites the moving entities are re-localized at every chunk
boundary and blended into a filtered track at $\alpha=0.6$, with re-detections
farther than 0.20\,m from the current track discarded as detector outliers.
Static entities keep their settle-time estimate.

\subsection{First-order hazard election}
\label{sec:ident}

The scene-level decision is made by a fixed rule, not by a model:
\begin{equation}
\begin{aligned}
\textsc{Hazard}(x) \coloneqq{}& \textsc{Protected}(x)\ \vee\ \textsc{Moving}(x)\ \vee \\
&\big(\lnot\textsc{Mentioned}(x)\wedge\lnot\textsc{Target}(x) \\
&\quad\wedge\ \textsc{NearPath}(x)\big).
\end{aligned}
\label{eq:fol}
\end{equation}
Table~\ref{tab:predicates} gives every predicate, its formal definition, and a
worked LIBERO-Safety instance, so we state here only the three properties of the
rule that the table cannot express.

First, the disjuncts are asymmetric by design. $\textsc{Protected}$ and
$\textsc{Moving}$ elect unconditionally, and in particular they are not
cancelled by mention, whereas the third disjunct is a conjunction that any one
of its literals can defeat. The consequence is that semantics may raise concern
but never remove protection, which is the fail-safe direction. Second,
$\textsc{NearPath}$ is a ranking at election time and a gate only at the
per-replan re-check. Election returns an ordered guard set rather than a single
argmax, protected and moving entities first and the remainder by distance to the
sanctioned reach segment. Retaining the top two is what converts identification
error from a fail-unsafe event into a recoverable one: under inverted, dropped,
and intact property priors the true hazard remained in the top two in all 19
probe scenes, against 3 of 19 for the argmax rule, because path geometry alone
keeps it in the set and the property ordering only permutes within it.

Third, it is worth stating what the rule is not. It is not synthesized per
scene, not retrieved, and not conditioned on a benchmark-specific table. On the
affordance suite the third disjunct's geometric literal is replaced by the
property literal $\textsc{Prop}$ of Table~\ref{tab:predicates}, and this
substitution is declared per suite in advance rather than selected per episode.
The rule is deliberately small enough to read, and every literal is separately
falsifiable, which is what makes the single-clause ablation of
Sec.~\ref{sec:ablation} possible at all.

\subsection{Grounding the predicates}
\label{sec:ground}

\textbf{What the model is asked.} The vision-language model is never asked which
object is dangerous, never asked to rank hazardousness, and never asked to
evaluate Eq.~\eqref{eq:fol}. It is asked only for the value of a named
predicate on a named object. Two probes motivate this restriction. A critic
probe that let the model compose the rule itself failed unsafe on three to eight
scenes, and a scene-wide hazard-scoring probe ranked a ketchup bottle above a
toy car. Both failures are decisions, and decisions are what the symbolic layer
retains.

\textbf{Reference and protection.} $\textsc{Mentioned}$ and $\textsc{Protected}$
are grounded per scene by a single structured query issued to
\texttt{claude-haiku-4-5} at temperature 0.3, repeated for three votes with a
two-of-three majority. The prompt supplies the instruction and the internal
object-name list, and requires JSON with two fields, the names the instruction
refers to including paraphrases and synonyms, and the names that are a human
body part, attached to a person, or a living being. The query is textual rather
than visual, which is deliberate: reference resolution and body-part
classification are language questions, and routing them through an image would
add a failure mode without adding information the name list does not already
carry. Because instructions are fixed per task, the majority-voted sets are
cached per scene and the runtime issues no API call at all.

\textbf{Properties.} $\textsc{Hot}$, $\textsc{Sharp}$, and $\textsc{Fragile}$
are grounded by an offline table built once over base object classes with five
votes and a three-of-five majority. Base-class normalization strips instance
indices, so an unseen instance of a known class inherits the class answer and an
unseen class is simply false on every property, which is the abstaining
direction. $\textsc{Hot}$ is then made state-dependent at runtime, since an
object of a hot class is an effective hazard only when a heat source lies within
0.15\,m of it, and a frying pan away from the stove is therefore not elected.

\textbf{Two interchangeable backends.} Every predicate marked in
Table~\ref{tab:predicates} can alternatively be grounded by a string heuristic
over head nouns and a fixed property vocabulary. On all 59 scenes of the two
benchmarks the heuristic and vision-language backends produce identical
elections, which is the strongest available evidence that the decision lives in
the logic rather than in the model. The vision-language backend nonetheless
resolves references that no string heuristic can, including the instruction
\textit{``bring it for me''}, under which the hand becomes $\textsc{Protected}$
without ever being named, and a plate held in a hand, which is correctly both
$\textsc{Mentioned}$ and $\textsc{Protected}$. The design finding may be stated
in one sentence: the vision-language model answers local predicates, and the
logic makes the decision.

\subsection{Three-regime routing}
\label{sec:routing}

Not every hazard is a keep-out region, and treating every hazard as one is a
failure mode that we measured. The same predicate layer routes each episode into
one of three regimes. In \textbf{Avoid}, a geometric hazard lies off the task
path and is enforced as described below. In \textbf{Disengage}, the hazard is
itself the destination, indicated by
$\textsc{Protected}(x)\wedge\textsc{Mentioned}(x)$ or by the guard overlapping
the parsed destination within 0.15\,m; the geometric keep-out is then disengaged
for the episode and safety is obtained from routing instead. This regime was
required by the human-handover suites, in which any keep-out region destroys the
task (Sec.~\ref{sec:ls}). In \textbf{Refuse}, the instruction itself is the
hazard, and a rule-based semantic judge over triples of verb, patient, and
hazard property refuses execution, with offline F1 0.968, precision 0.938, and
recall 1.000, measured in-domain as discussed in Sec.~\ref{sec:limitations}.
Routing and identification read the same predicates, so this is one rule system
rather than three.

\subsection{Enforcement}
\label{sec:enforce}

\textbf{From elected hazard to barrier.} Each elected guard contributes a
sphere-like keep-out region centered on its estimated position with effective
radius $r_{\mathrm{eff}}$ equal to its extent plus the end-effector radius plus a
safety pad. Reading an action chunk as a single-integrator rollout of
translational deltas gives predicted end-effector positions $p_j$ over the
horizon, and the barrier at step $j$ is
\begin{equation}
B_j(p)=\min_{m,q}\,\lVert p_j+q-o_m\rVert-r_m(j),
\end{equation}
minimized over guards $m$ and companion offsets $q$ that place additional
witness points at the wrist and fingertips, with $r_m(j)$ inflated linearly in
$j$ to bound obstacle motion within a chunk. A prefix is safe when it satisfies
the discrete control barrier chain $B_j\ge(1-\gamma)B_{j-1}$ with $\gamma=0.9$,
whose geometric decay is standard~\cite{agrawal2017dcbf} and for which we claim
no novelty. A semantic corridor exemption relaxes, and never removes, the margin
for rollout points inside a cylinder around the sanctioned approach segments
from the end-effector to the target and from the target to the destination. The
exemption is what allows a protective envelope wide enough to matter to coexist
with grasping and placing next to the guarded object, which a uniform envelope
cannot, since a uniform envelope cannot distinguish sanctioned approach from
hazard approach.

\textbf{How the barrier enters the steering.} Let $F(z,\mathrm{obs})$ denote the
frozen ten-step Euler denoise map and let $M(a_{1:L})$ denote the minimum chain
margin of the executed $L$-step prefix. All operators search the noise preimage
$S=\{z:M(F(z))\ge0\}$, so that the emitted action is $F(z^\star)$ for some seed
$z^\star$. \textsc{Select} draws $K{=}8$ seeds that share one vision-language
key-value cache and scores the decoded candidates lexicographically by
feasibility, then margin, then progress toward the target. \textsc{Tilt}
reweights and resamples the particle population at intermediate denoising steps
under soft margin potentials, at essentially no additional cost. \textsc{Ascend}
takes reverse-mode gradient steps of the prefix margin with respect to $z$
through all ten Euler steps at 57\,ms per step, with the backbone untouched. A
soft annealed repulsor field on the velocity target complements these operators
and is the only mechanism we found that navigates obstacle-on-grasp geometry, in
which the target sits inside the guarded region and every hard projection locks
up. The correction schedule trusts the flow early, corrects softly at
intermediate steps, and enforces exactly at the final step, since always-hard
repair from pure noise is a documented failure pattern for in-denoising safety.

\textbf{Why steering yields better actions than repair.} Three properties follow
from searching over $z$ rather than editing the decoded action. The emitted
chunk is a sample the frozen policy would itself have produced, so it lies on
the policy's own action manifold and inherits its temporal coherence, whereas
projecting a decoded chunk onto a constraint set produces trajectories the
policy never generates and cannot continue from. The corridor exemption and the
soft repulsor act on the flow rather than against it, which is what keeps the
system moving in the geometries where hard projection deadlocks. Finally, the
search is free to change the plan rather than only the increment, which is the
difference between routing around a hazard and stalling in front of it.

\textbf{Certificate.} The acceptance stage re-evaluates the chain on the
executed prefix of the final decoded chunk by explicit rollout of the actuation
model.

\begin{proposition}[Checked prefix certificate]\label{prop:cert}
If the acceptance check passes, the executed prefix satisfies the chain
regardless of how the chunk was produced. If all $K$ candidates fail, the
candidate with the largest margin executes and its true margin is reported, so
that the failure is visible rather than silent.
\end{proposition}

Because soundness is a predicate of the emitted actions rather than of the
search that produced them, the certificate covers \textsc{Select},
\textsc{Tilt}, \textsc{Ascend}, and the repulsor identically. Correction-based
arguments, of the form that the repair enforces the chain, cannot establish
this, and prior in-denoising work does not re-verify its intermediate
corrections on the final chunk at all. Several non-claims should be stated
explicitly. The rotation channels are uncertified, tracking error is absorbed
into the safety margin rather than proved, and we claim no formal
perception-robustness radius, quantifying it empirically instead
(Sec.~\ref{sec:e3}).
